\label{sec:competition}

\subsection{Two-Party Presidential Vote as the Measure}

Electoral competition can be measured many ways: incumbent vote share,
margin of victory in general elections, challenger quality, or campaign
spending.  For comparisons across algorithmic and enacted maps, we need a
measure that is (a) available for every district, (b) exogenous to the
specific candidates running, and (c) comparable across redistricting cycles.

Presidential vote share satisfies all three criteria.  Presidential election
results are available at the precinct level for the 2020 election and can be
reaggregated to any district geometry~\citep{deford2021redistricting}.
Presidential vote is less sensitive to candidate-specific idiosyncrasies than
congressional vote; it represents the underlying partisan lean of the
electorate rather than the personal vote of a specific incumbent.  It is
standard in the redistricting literature for precisely these
reasons~\citep{gelman1994unified}.

We compute the two-party Democratic presidential vote share $D_i$ for each
district $i$ as:
\begin{equation}
  D_i = \frac{\text{Biden}_i}{\text{Biden}_i + \text{Trump}_i}
\end{equation}
where $\text{Biden}_i$ and $\text{Trump}_i$ denote votes cast in 2020 in
district $i$'s territory.

\subsection{Competitive and Swing Thresholds}

We adopt two competition thresholds:

\begin{definition}[Competitive District]
A district is \emph{competitive} if $|D_i - 0.5| < 0.10$, i.e., if the
two-party Democratic presidential vote share falls between 40\% and 60\%.
\end{definition}

\begin{definition}[Swing District]
A district is a \emph{swing district} if $|D_i - 0.5| < 0.05$, i.e., if
the two-party Democratic presidential vote share falls between 45\% and 55\%.
\end{definition}

The 10-point competitive threshold is widely used in the congressional
elections literature and corresponds to a district that could plausibly flip
parties with a national wave of 5--8 points~\citep{mcdonald2006new}.  The
5-point swing threshold identifies the most hotly contested districts where
candidate quality and campaign resources plausibly determine the outcome.

\subsection{Partisan Lean Classification}

Among competitive districts, we classify partisan lean as follows:
\begin{itemize}
  \item \textbf{Lean Democratic:} $0.50 \leq D_i < 0.60$
  \item \textbf{Lean Republican:} $0.40 < D_i < 0.50$
  \item \textbf{True Toss-up:} $0.475 \leq D_i \leq 0.525$
\end{itemize}
This allows us to assess whether competitive algorithmic districts are
systematically advantaged for one party.

\subsection{Limitations of the Presidential Vote Proxy}

The presidential vote proxy has well-known limitations.  First, it measures
the 2020 election's specific partisan environment; if the 2020 election was
unusually nationalized, the proxy may overstate the true competitiveness of
districts whose local political dynamics differ from the national trend.
Second, incumbency effects---which can shift congressional vote share by 5--10
points relative to presidential results~\citep{ansolabehere2000incumbency}---mean
that a ``competitive'' district by presidential vote may in practice be safe
for an incumbent.  We address this limitation in Section~6 (longitudinal
projection) by examining stability under vote-share drift.
